\chapter{Codici MATLAB: Controllo Predittivo MPC}

In questa appendice vengono riportati i codici MATLAB utilizzati. Il documento è strutturato per essere facilmente navigabile: partendo dallo script principale, è possibile esplorare tutte le funzioni richiamate.

\section*{Script Principale MPC}
\subsection*{File: \texttt{MPC.m}}
\label{sec:MPC}

Script MATLAB principale per la simulazione MPC dell'F-16 (Configurazione, Inizializzazione, Ciclo Closed-Loop). In questo script vengono richiamate le seguenti funzioni, definite nelle sezioni successive:
\begin{itemize}
    \item \hyperref[sec:startup]{\texttt{startup.m}}
    \item \hyperref[sec:funzioni_mpc_discretizza_modello]{\texttt{funzioni\_mpc/discretizza\_modello.m}}
    \item \hyperref[sec:funzioni_mpc_imposta_vincoli]{\texttt{funzioni\_mpc/imposta\_vincoli.m}}
    \item \hyperref[sec:funzioni_mpc_setup_mpc]{\texttt{funzioni\_mpc/setup\_mpc.m}}
    \item \hyperref[sec:funzioni_mpc_simula_mpc]{\texttt{funzioni\_mpc/simula\_mpc.m}}
    \item \hyperref[sec:funzioni_mpc_cis]{\texttt{funzioni\_mpc/cis.m}}
    \item \hyperref[sec:funzioni_mpc_plot_cis]{\texttt{funzioni\_mpc/plot\_cis.m}}
    \item \hyperref[sec:funzioni_mpc_plot_mpc_cost_3d]{\texttt{funzioni\_mpc/plot\_mpc\_cost\_3d.m}}
    \item \hyperref[sec:funzioni_mpc_plot_risultati]{\texttt{funzioni\_mpc/plot\_risultati.m}}
\end{itemize}

\begin{lstlisting}[language=Matlab, caption={Script Principale MPC}, label={lst:MPC}]
% =========================================================================
% Tesi Triennale - MPC Longitudinale F-16
% Simulazione con Conversione Continuo-Discreto (c2d) e Radianti
% Ing. Leggeri Leonardo
% =========================================================================

% Inizializzazione Path
% Aggiungiamo prima la cartella locale Startup in cima al path per evitare
% conflitti con startup_project di altri progetti (es. Tre Camere)
addpath(fullfile(pwd, 'Startup'), '-begin');
startup_project();

% =========================================================================
% TABS AUTOMATICI: Inserisce tutti i grafici in un unico pannello a schede
% =========================================================================
set(0, 'DefaultFigureWindowStyle', 'docked');

%% 1. Definizione del Sistema (TEMPO CONTINUO)
nx = 4; % Stati: [theta,q,U,W]'
nu = 3; % Ingressi: [T, dele, dlef]'
Ts = 0.1; % Tempo di campionamento in secondi

disp('Conversione del modello da Continuo a Discreto...');
[A_long_ds, B_ctrl_ds] = discretizza_modello(A_long, B_ctrl, nx, nu, Ts);

%% 2. Progetto LQR tramite funzione dedicata
% =========================================================================
% NOTA SUI PESI Q e R:
% Se vuoi modificare l'aggressività del controllore (pesi Q e R),
% apri il file "funzioni_lqr/progetta_LQR_discreto.m" e modificali lì dentro
% =========================================================================
disp('Progetto LQR discreto tramite funzione dedicata...');
[K, P, Q, R, A_cl] = progetta_LQR_discreto(A_long_ds, B_ctrl_ds);

    % I vincoli FISICI assoluti dell'F-16:
    % Thrust: [1000, 19000] lbf -> Newton
    % Elevator: [-25, 25] gradi
    % Flaps (LEF): [0, 25] gradi
    lbf2N = 4.44822;
    ft2m = 0.3048;
    
    U_min_abs = [1000 * lbf2N; -25 * pi/180; 0];  
    U_max_abs = [19000 * lbf2N; 25 * pi/180; 25 * pi/180]; 

    % Poiché l'MPC ragiona in "perturbazioni" (u_err), i vincoli su u_err 
    % devono essere sottratti del punto di trim (best_u)
    U_min = U_min_abs - [best_u(1); best_u(2); best_u(3)];
    U_max = U_max_abs - [best_u(1); best_u(2); best_u(3)];

    Rate_max = [10000 * lbf2N; 60 * pi/180; 25 * pi/180];

    % Stati: theta(rad), q(rad/s), u(m/s), w(m/s)
    X_max = [ 45 * pi/180;  60 * pi/180;  100 * ft2m;  85 * ft2m]; 
    X_min = [-45 * pi/180; -60 * pi/180; -100 * ft2m; -85 * ft2m];

[Fx, fx, Fu, fu, dU_min, dU_max, Gx, gx] = imposta_vincoli(X_min, X_max, U_min, U_max, Rate_max, Ts);

% Riferimento di Default (Origine del sistema linearizzato)
x_ref = zeros(nx, 1);
u_ref = zeros(nu, 1);

% =========================================================================
% ESEMPIO con target non nullo
% Affinché l'MPC funzioni correttamente sul modello linearizzato, il nuovo 
% target deve essere un punto di equilibrio valido: x_ref = A_ds*x_ref + B_ds*u_ref
%
%u_ref = [100; 2 * pi/180; 0]; % Es: +1000 lbs di spinta e 2 gradi di equilibratore
 %x_ref = (eye(nx) - A_long_ds) \ (B_ctrl_ds * u_ref); % Stato di equilibrio esatto associato
% =========================================================================

%% 4. Calcolo Control Invariant Set e Plot
disp('--- CALCOLO del Control Invariant Set ---');
[G_inf, g_inf] = cis(A_long_ds, B_ctrl_ds, x_ref, u_ref, Fx, fx, Fu, fu, Q, R);

%% 5. Setup Problema MPC 
N = 40; % Orizzonte predittivo 
mpc_prob = setup_mpc(N, nx, nu, A_long_ds, B_ctrl_ds, Q, P, R, U_min, U_max, Gx, gx, G_inf, g_inf, x_ref, u_ref);

%% 4b. Plot 3D dei Set Invarianti e Feasible Set N-Step
% Inizializza Dashboard
tabs = crea_dashboard_mpc();
plot_cis(G_inf, g_inf, x_ref, mpc_prob, A_long_ds, dU_max, tabs.cis_3d);

%% 6. Simulazione MPC Completa 
disp('--- Avvio Ottimizzazione e Simulazione MPC ---');
% Ordine: [theta; q; u; w]
x_iniziale = [25 * (pi/180);  % theta: +25 gradi (cabrata)
              3 * (pi/180);  % q: +20 gradi/s (velocità di beccheggio)
              2;             % u: +20 ft/s (deviazione velocità forward)
              5];            % w: +20 ft/s (deviazione velocità verticale)
t_sim = 100; % tempo di simulazione

[storia_x, storia_u, storia_costo] = simula_mpc(mpc_prob, x_iniziale, t_sim, A_long_ds, B_ctrl_ds, dU_max);
disp('Ottimizzazione Riuscita. Il modello è matematicamente solido.');

disp('---------------------------------------------------');
disp('VERIFICA RAGGIUNGIMENTO TARGET (Ultimo Step)');
disp('Stato Finale Raggiunto (storia_x(:,end)):');
disp(storia_x(:,end));
disp('Target Desiderato (x_ref):');
disp(x_ref);
disp('Errore Assoluto [theta; q; U; W]:');
disp(abs(storia_x(:,end) - x_ref));
disp('---------------------------------------------------');

%% 7. Grafici 
plot_risultati(t_sim, storia_x, storia_u, U_min, U_max, x_ref, u_ref, tabs.stati, tabs.attuatori);

%% 8. Plot Funzionale di Costo MPC 3D
plot_mpc_cost_3d(mpc_prob, A_long_ds, dU_max, storia_x, storia_costo, Ts, tabs.cost_3d);
\end{lstlisting}

\section*{Inizializzazione Workspace}
\subsection*{File: \texttt{startup.m}}
\label{sec:startup}

Script di inizializzazione dell'ambiente di lavoro e configurazione automatica dei percorsi di progetto.

\begin{lstlisting}[language=Matlab, caption={Inizializzazione Workspace}, label={lst:startup}]
% =========================================================================
% Copyright (c) 2026 Ing. Leggeri Leonardo
% Tutti i diritti riservati.
%
% ATTENZIONE: Questo software e il relativo codice sorgente sono di proprietà 
% esclusiva dell'Ing. Leggeri Leonardo. È severamente vietata la copia, 
% la distribuzione, la modifica o la vendita a terzi senza l'esplicito 
% consenso scritto dell'autore. La distribuzione o la vendita non autorizzata 
% costituisce reato ed è perseguibile penalmente secondo le leggi vigenti.
% =========================================================================

% STARTUP.m
% Questo script viene eseguito automaticamente da MATLAB all'avvio 
% (se la cartella corrente è impostata su questa directory)
% oppure all'apertura del Progetto MATLAB.

disp('--- Esecuzione Automatica di Startup ---');
addpath(fullfile(pwd, 'Startup'));
startup_project();
\end{lstlisting}

\section*{Discretizzazione del Modello}
\subsection*{File: \texttt{funzioni\_mpc/discretizza\_modello.m}}
\label{sec:funzioni_mpc_discretizza_modello}

Funzione per la discretizzazione del modello continuo dello spazio di stato (c2d).

\begin{lstlisting}[language=Matlab, caption={Discretizzazione del Modello}, label={lst:funzioni_mpc_discretizza_modello}]
% =========================================================================
% Copyright (c) 2026 Ing. Leggeri Leonardo
% Tutti i diritti riservati.
%
% ATTENZIONE: Questo software e il relativo codice sorgente sono di proprietà 
% esclusiva dell'Ing. Leggeri Leonardo. È severamente vietata la copia, 
% la distribuzione, la modifica o la vendita a terzi senza l'esplicito 
% consenso scritto dell'autore. La distribuzione o la vendita non autorizzata 
% costituisce reato ed è perseguibile penalmente secondo le leggi vigenti.
% =========================================================================

function [A_long_ds, B_ctrl_ds] = discretizza_modello(A_long, B_ctrl, nx, nu, Ts)
    % Converte il sistema da tempo continuo a discreto
    sys_c = ss(A_long, B_ctrl, eye(nx), zeros(nx, nu));
    sys_d = c2d(sys_c, Ts, 'zoh'); % Zero-Order Hold
    A_long_ds = sys_d.A;
    B_ctrl_ds = sys_d.B;
end
\end{lstlisting}

\section*{Definizione dei Vincoli (Constraints)}
\subsection*{File: \texttt{funzioni\_mpc/imposta\_vincoli.m}}
\label{sec:funzioni_mpc_imposta_vincoli}

Funzione per la definizione rigorosa dei vincoli fisici e operativi (su stati e ingressi).

\begin{lstlisting}[language=Matlab, caption={Definizione dei Vincoli (Constraints)}, label={lst:funzioni_mpc_imposta_vincoli}]
% =========================================================================
% Copyright (c) 2026 Ing. Leggeri Leonardo
% Tutti i diritti riservati.
%
% ATTENZIONE: Questo software e il relativo codice sorgente sono di proprietà 
% esclusiva dell'Ing. Leggeri Leonardo. È severamente vietata la copia, 
% la distribuzione, la modifica o la vendita a terzi senza l'esplicito 
% consenso scritto dell'autore. La distribuzione o la vendita non autorizzata 
% costituisce reato ed è perseguibile penalmente secondo le leggi vigenti.
% =========================================================================

function [Fx, fx, Fu, fu, dU_min, dU_max, Gx, gx] = imposta_vincoli(X_min, X_max, U_min, U_max, Rate_max, Ts)
    % IMPOSTA_VINCOLI Costruisce le matrici poliedriche dei vincoli per l'MPC.
    % È una funzione generalizzata per qualsiasi sistema.
    %
    % Input:
    %   X_min, X_max: Limiti inferiori e superiori degli stati (nx x 1)
    %   U_min, U_max: Limiti inferiori e superiori degli ingressi (nu x 1)
    %   Rate_max: Massimo rateo di variazione degli ingressi (nu x 1)
    %   Ts: Tempo di campionamento
    %
    % Output:
    %   Fx, fx: Matrici vincoli di stato Fx * x <= fx
    %   Fu, fu: Matrici vincoli di ingresso Fu * u <= fu
    %   dU_min, dU_max: Vincoli di rateo
    %   Gx, gx: Copia di Fx, fx (per retrocompatibilità)

    nx = length(X_min);
    nu = length(U_min);

    % Vincoli di Rateo
    dU_max = Rate_max * Ts; 
    dU_min = -dU_max;

    % Vincoli Poliedrici di Stato
    Fx = [eye(nx); -eye(nx)]; 
    fx = [X_max; -X_min];

    % Vincoli Poliedrici di Ingresso
    Fu = [eye(nu); -eye(nu)]; 
    fu = [U_max; -U_min];
    
    % Per retrocompatibilità con altri codici MPC pre-esistenti
    Gx = Fx; 
    gx = fx; 
end
\end{lstlisting}

\section*{Setup e Configurazione MPC}
\subsection*{File: \texttt{funzioni\_mpc/setup\_mpc.m}}
\label{sec:funzioni_mpc_setup_mpc}

Funzione per la sintesi e l'impostazione del controllore predittivo (pesi Q e R, orizzonte predittivo, vincoli terminali).

\begin{lstlisting}[language=Matlab, caption={Setup e Configurazione MPC}, label={lst:funzioni_mpc_setup_mpc}]
% =========================================================================
% Copyright (c) 2026 Ing. Leggeri Leonardo
% Tutti i diritti riservati.
%
% ATTENZIONE: Questo software e il relativo codice sorgente sono di proprietà 
% esclusiva dell'Ing. Leggeri Leonardo. È severamente vietata la copia, 
% la distribuzione, la modifica o la vendita a terzi senza l'esplicito 
% consenso scritto dell'autore. La distribuzione o la vendita non autorizzata 
% costituisce reato ed è perseguibile penalmente secondo le leggi vigenti.
% =========================================================================

function mpc_prob = setup_mpc(N, nx, nu, A_long_ds, B_ctrl_ds, Q, P, R, U_min, U_max, Gx, gx, G_inf, g_inf, x_ref, u_ref)
    % SETUP_MPC Prepara tutte le matrici per il quadprog
    n_vars = N*nu + N*nx; 
    
    R_blk = kron(eye(N), R);
    Q_blk = blkdiag(kron(eye(N-1), Q), P);
    H = 2 * blkdiag(R_blk, Q_blk);         
    
    % Termine lineare della funzione di costo per l'inseguimento del riferimento
    f = zeros(n_vars, 1);                  
    for k = 1:N
        f((k-1)*nu + 1 : k*nu) = -2 * R * u_ref;
    end
    for k = 1:N-1
        f(N*nu + (k-1)*nx + 1 : N*nu + k*nx) = -2 * Q * x_ref;
    end
    f(N*nu + (N-1)*nx + 1 : N*nu + N*nx) = -2 * P * x_ref;
    
    Aeq_base = zeros(N*nx, n_vars);
    for k = 1:N
        Aeq_base((k-1)*nx + 1 : k*nx, (k-1)*nu + 1 : k*nu) = -B_ctrl_ds;
        Aeq_base((k-1)*nx + 1 : k*nx, N*nu + (k-1)*nx + 1 : N*nu + k*nx) = eye(nx);
        if k > 1
            Aeq_base((k-1)*nx + 1 : k*nx, N*nu + (k-2)*nx + 1 : N*nu + (k-1)*nx) = -A_long_ds;
        end
    end
    
    lb = -inf(n_vars, 1); ub = inf(n_vars, 1); 
    for k = 1:N
        lb((k-1)*nu + 1 : k*nu) = U_min;
        ub((k-1)*nu + 1 : k*nu) = U_max;
    end
    
    A_ineq_stat = []; b_ineq_stat = [];
    for j = 1:N-1
        A_t = zeros(size(Gx, 1), n_vars);
        A_t(:, N*nu + (j-1)*nx + 1 : N*nu + j*nx) = Gx;
        A_ineq_stat = [A_ineq_stat; A_t];
        b_ineq_stat = [b_ineq_stat; gx];
    end
    
    A_t = zeros(size(G_inf, 1), n_vars);
    A_t(:, N*nu + (N-1)*nx + 1 : N*nu + N*nx) = G_inf;
    A_ineq_stat = [A_ineq_stat; A_t];
    b_ineq_stat = [b_ineq_stat; g_inf];
    
    % Salvataggio in struttura
    mpc_prob.H = H;
    mpc_prob.f = f;
    mpc_prob.Aeq_base = Aeq_base;
    mpc_prob.lb = lb;
    mpc_prob.ub = ub;
    mpc_prob.A_ineq_stat = A_ineq_stat;
    mpc_prob.b_ineq_stat = b_ineq_stat;
    mpc_prob.n_vars = n_vars;
    mpc_prob.N = N;
    mpc_prob.nx = nx;
    mpc_prob.nu = nu;
end
\end{lstlisting}

\section*{Simulatore MPC Closed-Loop}
\subsection*{File: \texttt{funzioni\_mpc/simula\_mpc.m}}
\label{sec:funzioni_mpc_simula_mpc}

Funzione core per la simulazione temporale e l'evoluzione del ciclo chiuso con controllo MPC.

\begin{lstlisting}[language=Matlab, caption={Simulatore MPC Closed-Loop}, label={lst:funzioni_mpc_simula_mpc}]
% =========================================================================
% Copyright (c) 2026 Ing. Leggeri Leonardo
% Tutti i diritti riservati.
%
% ATTENZIONE: Questo software e il relativo codice sorgente sono di proprietà 
% esclusiva dell'Ing. Leggeri Leonardo. È severamente vietata la copia, 
% la distribuzione, la modifica o la vendita a terzi senza l'esplicito 
% consenso scritto dell'autore. La distribuzione o la vendita non autorizzata 
% costituisce reato ed è perseguibile penalmente secondo le leggi vigenti.
% =========================================================================

function [storia_x, storia_u, storia_costo] = simula_mpc(mpc_prob, x_iniziale, t_sim, A_long_ds, B_ctrl_ds, dU_max)
    % SIMULA_MPC Risolve il problema quadratico e simula l'evoluzione del sistema
    nx = mpc_prob.nx;
    nu = mpc_prob.nu;
    N = mpc_prob.N;
    n_vars = mpc_prob.n_vars;
    
    storia_x = zeros(nx, t_sim+1); storia_x(:,1) = x_iniziale;
    storia_u = zeros(nu, t_sim);
    storia_costo = zeros(1, t_sim);
    u_previous = [0;0;0]; 
    options = optimoptions('quadprog', 'Display', 'off');
    % --- Precomputazione Matrici Invarianti ---
    A_rate = zeros(nu*2*N, n_vars); 
    b_rate_base = zeros(nu*2*N, 1);
    for k = 1:N
        idx_u = (k-1)*nu+1:k*nu;
        A_rate((k-1)*2*nu+1:k*2*nu, idx_u) = [eye(nu); -eye(nu)];
        if k > 1
            idx_u_prev = (k-2)*nu+1:(k-1)*nu;
            A_rate((k-1)*2*nu+1:k*2*nu, idx_u_prev) = [-eye(nu); eye(nu)];
            b_rate_base((k-1)*2*nu+1:k*2*nu) = [dU_max; dU_max];
        end
    end
    A_ineq_tot = [mpc_prob.A_ineq_stat; A_rate];
    % ------------------------------------------

    for t = 1:t_sim
        beq = zeros(N*nx, 1);
        beq(1:nx) = A_long_ds * storia_x(:,t); 
        
        b_rate = b_rate_base;
        b_rate(1:2*nu) = [dU_max + u_previous; dU_max - u_previous];
        b_ineq_tot = [mpc_prob.b_ineq_stat; b_rate];
        [z_opt, fval, exitflag] = quadprog(mpc_prob.H, mpc_prob.f, A_ineq_tot, b_ineq_tot, mpc_prob.Aeq_base, beq, mpc_prob.lb, mpc_prob.ub, [], options);
        if exitflag < 0
            error('Infeasible! Il punto allo step %d è fuori da X_N. Riduci leggermente la severità di x_iniziale.', t);
        end
        
        storia_costo(t) = fval;
        
        % L'animazione live è stata rimossa per massimizzare le prestazioni.
        
        u_applicata = z_opt(1:nu);
        storia_u(:, t) = u_applicata;
        
        storia_x(:, t+1) = A_long_ds * storia_x(:,t) + B_ctrl_ds * u_applicata;
        u_previous = u_applicata;
    end
    
    % Plot della traiettoria completa tutto in una volta alla fine
    plot3(storia_x(3,:), storia_x(4,:), storia_x(2,:), '-b', 'LineWidth', 2.5, 'DisplayName', 'Traiettoria Effettiva');
    
    % Plot punto di partenza
    plot3(x_iniziale(3), x_iniziale(4), x_iniziale(2), 'k*', 'MarkerSize', 10, 'LineWidth', 2, 'DisplayName', 'Partenza');
    text(x_iniziale(3), x_iniziale(4), x_iniziale(2)+0.05, ' Partenza', 'FontWeight', 'bold');
    
    % Aggiungi un pallino rosso per ogni passo effettivamente compiuto alla fine
    plot3(storia_x(3,:), storia_x(4,:), storia_x(2,:), 'ro', 'MarkerSize', 6, 'MarkerFaceColor', 'r', 'DisplayName', 'Passi MPC');
    
    % Plot punto di arrivo finale
    plot3(storia_x(3,end), storia_x(4,end), storia_x(2,end), 'kp', 'MarkerSize', 15, 'MarkerFaceColor', 'k', 'DisplayName', 'Arrivo Effettivo');
    text(storia_x(3,end), storia_x(4,end), storia_x(2,end)-0.05, ' Arrivo', 'FontWeight', 'bold', 'Color', 'k');
    
    % Aggiorna la legenda includendo i nuovi plot
    legend('show', 'Location', 'best', 'Interpreter', 'latex');
end
\end{lstlisting}

\section*{Calcolo Insieme Invariante (CIS)}
\subsection*{File: \texttt{funzioni\_mpc/cis.m}}
\label{sec:funzioni_mpc_cis}

Funzione per il calcolo del Control Invariant Set e del Maximum Control Invariant Set (Insieme Invariante di Controllo).

\begin{lstlisting}[language=Matlab, caption={Calcolo Insieme Invariante (CIS)}, label={lst:funzioni_mpc_cis}]
% =========================================================================
% Copyright (c) 2026 Ing. Leggeri Leonardo
% Tutti i diritti riservati.
%
% ATTENZIONE: Questo software e il relativo codice sorgente sono di proprietà 
% esclusiva dell'Ing. Leggeri Leonardo. È severamente vietata la copia, 
% la distribuzione, la modifica o la vendita a terzi senza l'esplicito 
% consenso scritto dell'autore. La distribuzione o la vendita non autorizzata 
% costituisce reato ed è perseguibile penalmente secondo le leggi vigenti.
% =========================================================================

function [G_inf, g_inf] = cis(A, B, x_ref, u_ref, Fx, fx, Fu, fu, Q, R)
    % Control Invariant Set Iterativo tramite Metodo Matriciale (No Polyhedron)
    [K, ~, ~] = dlqr(A, B, Q, R);
    A_cl = A - B*K;
    
    % Traslazione dei vincoli rispetto al riferimento z = x - x_ref
    Gx = Fx; 
    gx = fx - Fx*x_ref;
    
    Gu_x = [-K; K]; 
    gu_x = [fu(1:3) - Fu(1:3,:)*u_ref; fu(4:6) - Fu(4:6,:)*u_ref];
    
    G = [Gx; Gu_x]; 
    g = [gx; gu_x];

    G_inf = G; g_inf = g;
    max_iter = 100; tol = 1e-6;
    for i = 1:max_iter
        G_next = G * (A_cl^i);
        if norm(G_next, inf) < tol
            break;
        end
        G_inf = [G_inf; G_next];
        g_inf = [g_inf; g];
    end
    
    % Ritraslazione di g_inf in modo che G_inf * x <= g_inf 
    % (dato che attualmente G_inf * z <= g_inf)
    g_inf = g_inf + G_inf * x_ref;
end
\end{lstlisting}

\section*{Visualizzazione Insiemi Invarianti}
\subsection*{File: \texttt{funzioni\_mpc/plot\_cis.m}}
\label{sec:funzioni_mpc_plot_cis}

Funzione per la visualizzazione 3D e 2D (tramite proiezioni) degli insiemi invarianti calcolati.

\begin{lstlisting}[language=Matlab, caption={Visualizzazione Insiemi Invarianti}, label={lst:funzioni_mpc_plot_cis}]
% =========================================================================
% Copyright (c) 2026 Ing. Leggeri Leonardo
% Tutti i diritti riservati.
%
% ATTENZIONE: Questo software e il relativo codice sorgente sono di proprietà 
% esclusiva dell'Ing. Leggeri Leonardo. È severamente vietata la copia, 
% la distribuzione, la modifica o la vendita a terzi senza l'esplicito 
% consenso scritto dell'autore. La distribuzione o la vendita non autorizzata 
% costituisce reato ed è perseguibile penalmente secondo le leggi vigenti.
% =========================================================================

function plot_cis(G_inf, g_inf, x_ref, mpc_prob, A_long_ds, dU_max, parent_tab)
    if nargin < 3
        x_ref = zeros(4,1);
    end
    if nargin < 7
        fig1 = figure('Name', 'Control Invariant Set 3D', 'Color', 'w', 'Position', [50, 50, 900, 700]);
        parent_tab = fig1;
    end
    
    % PLOT_CIS Disegna in 3D il Control Invariant Set
    ax = axes('Parent', parent_tab);
    hold(ax, 'on'); grid(ax, 'on'); view(ax, 3); 
    title(ax, ['Poliedro $\mathcal{X}_f$ in 3D (Fetta $\theta = ', num2str(x_ref(1)), '$)'], 'Interpreter', 'latex', 'FontSize', 14);
    
    % Assi aggiornati al tuo vettore di stato: [theta, q, U, W]
    xlabel('u [m/s]'); ylabel('w [m/s]'); zlabel('q [rad/s]');
    
    % Griglia ad alta risoluzione centrata vicino all'origine per trovare il CIS
    [X_U, X_W, X_q] = meshgrid(linspace(-30, 30, 40), ... % Range per u (Stato 3)
                               linspace(-30, 30, 40), ... % Range per w (Stato 4)
                               linspace(-40 * pi/180, 40 * pi/180, 40)); % Range per q (Stato 2)
    
    X_U_f = X_U(:); X_W_f = X_W(:); X_q_f = X_q(:); 
    
    % Imposta la fetta in modo che combaci con la theta del target
    theta_slice = x_ref(1); 
    X_theta_f = theta_slice * ones(size(X_U_f)); % Fetta corrispondente al target
    
    Validi_inf = true(size(X_U_f));
    for j = 1:size(G_inf, 1)
        Valore = G_inf(j,1)*X_theta_f + G_inf(j,2)*X_q_f + G_inf(j,3)*X_U_f + G_inf(j,4)*X_W_f;
        Validi_inf = Validi_inf & (Valore <= g_inf(j));
    end
    
    PX_inf = X_U_f(Validi_inf); PY_inf = X_W_f(Validi_inf); PZ_inf = X_q_f(Validi_inf);
    
    if length(PX_inf) > 4
        K_hull_inf = convhull(PX_inf, PY_inf, PZ_inf);
        h_inf = trisurf(K_hull_inf, PX_inf, PY_inf, PZ_inf, 'Parent', ax, 'FaceColor', 'c', 'FaceAlpha', 0.8, 'EdgeColor', 'b', 'EdgeAlpha', 0.3);
        leg_handles = h_inf;
        leg_names = {'Control Invariant Set $\mathcal{O}_\infty$'};
    else
        disp('ATTENZIONE: Nessun punto valido trovato per il plot di O_inf sulla fetta scelta.');
        leg_handles = [];
        leg_names = {};
    end

    % --- NOVITÀ: Calcolo e Plot N-Step Feasible Set ---
    if nargin >= 6
        disp('Calcolo del Set Feasible N-Step X_N in corso (esplorazione vincoli)...');
        nx = mpc_prob.nx;
        nu = mpc_prob.nu;
        N = mpc_prob.N;
        n_vars = mpc_prob.n_vars;
        
        options = optimoptions('linprog', 'Display', 'off');
        A_rate = zeros(nu*2*N, n_vars); 
        b_rate_base = zeros(nu*2*N, 1);
        for k = 1:N
            idx_u = (k-1)*nu+1:k*nu;
            A_rate((k-1)*2*nu+1:k*2*nu, idx_u) = [eye(nu); -eye(nu)];
            if k > 1
                idx_u_prev = (k-2)*nu+1:(k-1)*nu;
                A_rate((k-1)*2*nu+1:k*2*nu, idx_u_prev) = [-eye(nu); eye(nu)];
                b_rate_base((k-1)*2*nu+1:k*2*nu) = [dU_max; dU_max];
            end
        end
        A_ineq_tot = [mpc_prob.A_ineq_stat; A_rate];
        
        u_previous = [0;0;0];
        b_rate = b_rate_base;
        b_rate(1:2*nu) = [dU_max + u_previous; dU_max - u_previous];
        b_ineq_tot = [mpc_prob.b_ineq_stat; b_rate];

        % Griglia a bassa risoluzione dedicata SOLO all'N-Step Set per velocizzare linprog
        [X_U_N, X_W_N, X_q_N] = meshgrid(linspace(-40, 40, 15), ... 
                                         linspace(-40, 40, 15), ... 
                                         linspace(-40 * pi/180, 40 * pi/180, 15));
        
        X_U_N_f = X_U_N(:); X_W_N_f = X_W_N(:); X_q_N_f = X_q_N(:); 
        f_lin = zeros(n_vars, 1);
        Validi_N = false(size(X_U_N_f));
        
        for i = 1:length(X_U_N_f)
            % Check rapido se è già in O_inf per saltare linprog
            x0_test = [theta_slice; X_q_N_f(i); X_U_N_f(i); X_W_N_f(i)];
            
            is_in_O_inf = true;
            for j = 1:size(G_inf, 1)
                if G_inf(j,:) * x0_test > g_inf(j)
                    is_in_O_inf = false;
                    break;
                end
            end
            
            if is_in_O_inf
                Validi_N(i) = true;
                continue;
            end
            
            % Se non è in O_inf, testa con linprog
            beq = zeros(N*nx, 1);
            beq(1:nx) = A_long_ds * x0_test; 
            
            [~, ~, exitflag] = linprog(f_lin, A_ineq_tot, b_ineq_tot, mpc_prob.Aeq_base, beq, mpc_prob.lb, mpc_prob.ub, options);
            if exitflag >= 0 || exitflag == -3
                Validi_N(i) = true;
            end
        end
        
        PX_N = X_U_N_f(Validi_N); PY_N = X_W_N_f(Validi_N); PZ_N = X_q_N_f(Validi_N);
        if length(PX_N) > 4
            K_hull_N = convhull(PX_N, PY_N, PZ_N);
            h_N = trisurf(K_hull_N, PX_N, PY_N, PZ_N, 'Parent', ax, 'FaceColor', 'y', 'FaceAlpha', 0.2, 'EdgeColor', 'y', 'EdgeAlpha', 0.1);
            leg_handles(end+1) = h_N;
            leg_names{end+1} = sprintf('Feasible N-Step Set $\\mathcal{X}_{%d}$', N);
            title(ax, ['Poliedri $\mathcal{O}_\infty$ e $\mathcal{X}_{', num2str(N), '}$ in 3D (Fetta $\theta = ', num2str(theta_slice), '$)'], 'Interpreter', 'latex', 'FontSize', 14);
        end
    end
    
    if ~isempty(leg_handles)
        legend(leg_handles, leg_names, 'Location', 'best', 'Interpreter', 'latex');
    end
end
\end{lstlisting}

\section*{Analisi Costo MPC 3D}
\subsection*{File: \texttt{funzioni\_mpc/plot\_mpc\_cost\_3d.m}}
\label{sec:funzioni_mpc_plot_mpc_cost_3d}

Funzione per il plotting tridimensionale della funzione di costo associata al regolatore MPC.

\begin{lstlisting}[language=Matlab, caption={Analisi Costo MPC 3D}, label={lst:funzioni_mpc_plot_mpc_cost_3d}]
% =========================================================================
% Copyright (c) 2026 Ing. Leggeri Leonardo
% Tutti i diritti riservati.
%
% ATTENZIONE: Questo software e il relativo codice sorgente sono di proprietà 
% esclusiva dell'Ing. Leggeri Leonardo. È severamente vietata la copia, 
% la distribuzione, la modifica o la vendita a terzi senza l'esplicito 
% consenso scritto dell'autore. La distribuzione o la vendita non autorizzata 
% costituisce reato ed è perseguibile penalmente secondo le leggi vigenti.
% =========================================================================

function plot_mpc_cost_3d(mpc_prob, A_long_ds, dU_max, storia_x, storia_costo, Ts, parent_tab)
    % =========================================================================
    % PLOT_MPC_COST_3D - Visualizza il Funzionale di Costo Ottimo J*(x) dell'MPC
    % =========================================================================

    if nargin < 6
        Ts = 0.1;
    end
    if nargin < 7
        fig1 = figure('Name', sprintf('Costo MPC 3D (Ts = %g s)', Ts), 'Color', 'w', 'Position', [150 150 850 650]);
        parent_tab = fig1;
    end
    w_lim = 40;  % m/s
    q_lim = 1.0; % rad/s

    % Usiamo una griglia 20x20 per un calcolo rapido
    [W_grid, Q_rad] = meshgrid(linspace(-w_lim, w_lim, 20), linspace(-q_lim, q_lim, 20));

    nx = mpc_prob.nx;
    nu = mpc_prob.nu;
    N = mpc_prob.N;
    n_vars = mpc_prob.n_vars;

    % Recuperiamo P_ds approssimato (o potresti passarlo da MPC.m)
    [~, P_ds, ~] = dlqr(A_long_ds, mpc_prob.Aeq_base(1:nx, 1:nu)*(-1), blkdiag(10,10,1,1), eye(nu)); 
    
    V_surf = zeros(size(W_grid));
    for i = 1:size(W_grid, 1)
        for j = 1:size(W_grid, 2)
            stato_surf = [0; Q_rad(i,j); 0; W_grid(i,j)];
            V_surf(i,j) = stato_surf' * P_ds * stato_surf;
        end
    end

    %% --- FIGURA 3D ---
    ax = axes('Parent', parent_tab);
    hold(ax, 'on'); grid(ax, 'on');
    
    h_surf = surf(ax, W_grid, Q_rad, V_surf, 'EdgeColor', 'none', 'FaceAlpha', 0.65);
    colormap(ax, jet);
    cb = colorbar(ax);
    ylabel(cb, 'Costo Ottimo MPC $J^*(x_k)$', 'Interpreter', 'latex', 'FontSize', 12);
    
    max_v_traj = max(storia_costo);
    if isempty(max_v_traj) || isnan(max_v_traj)
        max_v_traj = 1;
    end

    % Traiettoria
    N_steps = length(storia_costo);
    t_discrete = 1:N_steps;
    t_fine = linspace(1, N_steps, N_steps * 10);
    
    W_smooth = pchip(t_discrete, storia_x(4,1:N_steps), t_fine);
    Q_smooth = pchip(t_discrete, storia_x(2,1:N_steps), t_fine);
    V_smooth = pchip(t_discrete, storia_costo, t_fine);
    
    h_line = plot3(ax, W_smooth, Q_smooth, V_smooth, '-r', 'LineWidth', 2);

    % Punti discreti veri
    plot3(ax, storia_x(4,1:N_steps), storia_x(2,1:N_steps), storia_costo, 'o', ...
          'MarkerEdgeColor', 'r', 'MarkerFaceColor', 'w', 'MarkerSize', 4);
          
    % Start
    h_start = plot3(ax, storia_x(4,1), storia_x(2,1), storia_costo(1), 'o', ...
                    'MarkerFaceColor', 'r', 'MarkerEdgeColor', 'k', 'MarkerSize', 8);
    
    % End
    h_end = plot3(ax, storia_x(4,N_steps), storia_x(2,N_steps), storia_costo(N_steps), 's', ...
          'MarkerFaceColor', 'g', 'MarkerEdgeColor', 'k', 'MarkerSize', 8);

    title(ax, sprintf('\\textbf{Funzionale di Costo MPC $J^*(x_k)$ a Tempo Discreto ($T_s = %g$ s)}', Ts), 'Interpreter', 'latex', 'FontSize', 16);
    xlabel(ax, 'Velocit\`a Verticale $w$ [m/s]', 'Interpreter', 'latex', 'FontSize', 12);
    ylabel(ax, 'Velocit\`a di Beccheggio $q$ [rad/s]', 'Interpreter', 'latex', 'FontSize', 12);
    zlabel(ax, 'Costo $J^*(x_k)$', 'Interpreter', 'latex', 'FontSize', 12);
    
    legend(ax, [h_surf, h_line, h_start, h_end], ...
           {'Superficie $J^*(x_k)$', 'Evoluzione di Stato', 'Partenza ($k=0$)', 'Arrivo'}, ...
           'Interpreter', 'latex', 'FontSize', 12, 'Location', 'northeast');
    
    zlim(ax, [0, max(max(V_surf(:)), max_v_traj) * 1.2]); 
    xlim(ax, [-w_lim, w_lim]);
    ylim(ax, [-q_lim, q_lim]);
    view(ax, -35, 30);
end
\end{lstlisting}

\section*{Plot e Grafici dei Risultati}
\subsection*{File: \texttt{funzioni\_mpc/plot\_risultati.m}}
\label{sec:funzioni_mpc_plot_risultati}

Funzione per la rappresentazione grafica delle risposte del sistema, evoluzione degli stati e comandi degli attuatori.

\begin{lstlisting}[language=Matlab, caption={Plot e Grafici dei Risultati}, label={lst:funzioni_mpc_plot_risultati}]
% =========================================================================
% Copyright (c) 2026 Ing. Leggeri Leonardo
% Tutti i diritti riservati.
%
% ATTENZIONE: Questo software e il relativo codice sorgente sono di proprietà 
% esclusiva dell'Ing. Leggeri Leonardo. È severamente vietata la copia, 
% la distribuzione, la modifica o la vendita a terzi senza l'esplicito 
% consenso scritto dell'autore. La distribuzione o la vendita non autorizzata 
% costituisce reato ed è perseguibile penalmente secondo le leggi vigenti.
% =========================================================================

%funzione per la creazione dei grafici dopo esecuzione script MPC.m
function plot_risultati(t_sim, storia_x, storia_u, U_min, U_max, x_ref, u_ref, tab_stati, tab_attuatori)
    if nargin < 6
        x_ref = zeros(4,1);
        u_ref = zeros(3,1);
    end
    if nargin < 8
        fig1 = figure('Name', 'Risultati MPC: Stati del Velivolo', 'Color', 'w', 'Position', [100 100 900 600]);
        tab_stati = fig1;
        fig2 = figure('Name', 'Risultati MPC: Sforzo degli Attuatori', 'Color', 'w', 'Position', [150 150 700 800]);
        tab_attuatori = fig2;
    end

    % PLOT_RISULTATI Disegna i grafici dell'evoluzione di stati e attuatori
    t_layout_x = tiledlayout(tab_stati, 2, 2, 'Padding', 'compact', 'TileSpacing', 'compact');
    
    nexttile(t_layout_x); plot(0:t_sim, storia_x(1,:), '-b', 'LineWidth', 1.5); hold on; yline(x_ref(1), 'r--', 'LineWidth', 1.2); title('Pitch Angle (\theta) [rad]'); grid on; 
    nexttile(t_layout_x); plot(0:t_sim, storia_x(2,:), '-r', 'LineWidth', 1.5); hold on; yline(x_ref(2), 'r--', 'LineWidth', 1.2); title('Pitch Rate (q) [rad/s]'); grid on; 
    nexttile(t_layout_x); plot(0:t_sim, storia_x(3,:), '-m', 'LineWidth', 1.5); hold on; yline(x_ref(3), 'r--', 'LineWidth', 1.2); title('Velocità Forward (u) [m/s]'); grid on; 
    nexttile(t_layout_x); plot(0:t_sim, storia_x(4,:), '-c', 'LineWidth', 1.5); hold on; yline(x_ref(4), 'r--', 'LineWidth', 1.2); title('Velocità Verticale (w) [m/s]'); grid on;
    lgd = legend({'Traiettoria', 'Target (x_{ref})'}, 'Orientation', 'horizontal');
    lgd.Layout.Tile = 'north';

    t_layout_u = tiledlayout(tab_attuatori, 3, 1, 'Padding', 'compact', 'TileSpacing', 'compact');
    
    nexttile(t_layout_u); stairs(0:t_sim-1, storia_u(1,:), '-g', 'LineWidth', 1.5); hold on;
    yline(U_max(1), 'k--'); yline(U_min(1), 'k--'); yline(u_ref(1), 'r--', 'LineWidth', 1.2); title(' Spinta [N]'); grid on;
    
    nexttile(t_layout_u); stairs(0:t_sim-1, storia_u(2,:) * 180 / pi, '-g', 'LineWidth', 1.5); hold on;
    yline(U_max(2) * 180 / pi, 'k--'); yline(U_min(2) * 180 / pi, 'k--'); yline(u_ref(2) * 180 / pi, 'r--', 'LineWidth', 1.2);
    title(' Elevatore [deg]'); grid on;
    
    nexttile(t_layout_u); stairs(0:t_sim-1, storia_u(3,:) * 180 / pi, '-g', 'LineWidth', 1.5); hold on;
    yline(U_max(3) * 180 / pi, 'k--'); yline(U_min(3) * 180 / pi, 'k--'); yline(u_ref(3) * 180 / pi, 'r--', 'LineWidth', 1.2);
    title(' Flap [deg]'); grid on;
end
\end{lstlisting}

